%%%%%%%%%%%%%%%%%%%%%%%%%%%%%%%%%%%%%%%%%%%%%%%%%%%%%%%%%%%%%%%%%%%%%%%%%%
%% Draft submission to INFORMS Journal on Computing (IJOC).
%% Built from the official INFORMS-IJOC-Template (v. Feb 2026).
%%%%%%%%%%%%%%%%%%%%%%%%%%%%%%%%%%%%%%%%%%%%%%%%%%%%%%%%%%%%%%%%%%%%%%%%%%
\documentclass[ijoc,sglanonrev]{informs4}
\usepackage{eqndefns-left}
\RequirePackage{tgtermes}
\RequirePackage{newtxtext}
\RequirePackage{newtxmath}
\RequirePackage{bm}
\RequirePackage{endnotes}

\OneAndAHalfSpacedXII

\usepackage{algorithm}
\usepackage{algpseudocode}
\usepackage{amsfonts}

\usepackage{natbib}
 \bibpunct[, ]{(}{)}{,}{a}{}{,}%
 \def\bibfont{\small}%
 \def\bibsep{\smallskipamount}%
 \def\bibhang{24pt}%
 \def\newblock{\ }%
 \def\BIBand{and}%

\EquationsNumberedThrough
\TheoremsNumberedThrough
\ECRepeatTheorems

\MANUSCRIPTNO{}

% Convenience macros
\newcommand{\Rspread}{R}
\newcommand{\FF}{F}

\begin{document}

\RUNAUTHOR{Dang}

\RUNTITLE{Adaptive Influence Minimization by Deferred Blocking}

\TITLE{Adaptive Influence Minimization via Deferred Frontier Blocking: Theory, Algorithms, and Exact Validation on Bounded-Treewidth Graphs}

\ARTICLEAUTHORS{%
\AUTHOR{Quang-Vinh Dang}
\AFF{British University Vietnam, Hung Yen, Vietnam, \EMAIL{vinh.dq4@buv.edu.vn}}
}

\ABSTRACT{%
Influence minimization (IMIN) asks for a small set of vertices whose removal
most reduces the spread of a fixed, already-seeded cascade -- the natural
model for blocking misinformation after it has started, rather than
preventing it before it starts. Existing algorithms are one-shot: the whole
blocking budget is committed before any part of the cascade is observed. We
ask whether committing adaptively, as the cascade unfolds, can only help,
and by how much. We give DEFER, a wrapper that turns any one-shot IMIN
planner into an adaptive policy by deferring the commitment of each planned
block until the moment it is actually exposed to the live cascade, and by
committing early only when waiting is provably not worth it. We prove that
deferred commitment is free (it never increases the final spread relative to
the wrapped planner), that the push-down commitment rule is the myopic
optimum for an exposed node protecting cheaply-replaceable children, and
that the resulting adaptivity gap is unbounded in the maximum out-degree of
the graph -- so the benefit of adaptivity is real and not an artifact of a
particular instance. We complement DEFER with TWIG, an exact
treewidth-bounded dynamic program that computes the true (not
Monte-Carlo-sampled) expected spread of any candidate blocking set, together
with a fixed-parameter-tractable reformulation and a provable separation
showing that the direction in which a tree decomposition is processed is a
first-class algorithmic choice, not a tuning detail. TWIG turns greedy IMIN
into a verifiable procedure on graphs small enough to admit exact
computation and lets us measure, for the first time free of sampling noise,
how often and how badly greedy departs from the true optimum. Together the
two contributions give a rigorously justified adaptive blocking algorithm
and an exact tool for validating any one-shot blocking heuristic against
ground truth, closing a verification gap that Monte-Carlo-based influence
estimation has left open throughout the influence-minimization literature.
}%

\FUNDING{The author received no external funding for this research.}

\KEYWORDS{influence minimization, adaptive optimization, misinformation blocking, treewidth, dynamic programming, network science}

\maketitle
%%%%%%%%%%%%%%%%%%%%%%%%%%%%%%%%%%%%%%%%%%%%%%%%%%%%%%%%%%%%%%%%%%%%%%

\section{Introduction}\label{sec:intro}

Misinformation and epidemic-style contagions on networks share a common
downstream control problem once they have already started: given a fixed
set of active or seeded sources and a limited budget of vertices that can be
removed (accounts suspended, individuals quarantined, links cut), which
vertices should be chosen to minimize the eventual spread? This is the
\emph{influence minimization} (IMIN) problem \citep{xie2023,xie2024}, the
natural complement of the much more heavily studied influence
\emph{maximization} problem \citep{kempe2003}: instead of choosing seeds to
spread a cascade, IMIN chooses a small blocking set to contain one. Unlike
maximization, the objective is not supermodular in general, so greedy
selection carries no approximation guarantee, and the problem is NP-hard and
hard to approximate \citep{xie2023}.

Every published IMIN algorithm -- AdvancedGreedy and GreedyReplace
\citep{xie2023,xie2024}, and the Sandwich lower/upper-bound family -- is
\emph{one-shot}: the entire budget is spent before the cascade is observed
even once. In practice a moderator or public-health authority does not have
to decide everything up front. They observe who has already been exposed,
round by round, and can adapt. This raises two questions that, to our
knowledge, have not been formally answered for IMIN: (i) can adapting to the
unfolding cascade only help, never hurt, relative to any one-shot plan; and
(ii) how large can the benefit of adaptivity be, in the worst case? A third,
orthogonal problem is one of verification: every existing IMIN algorithm is
evaluated by Monte-Carlo simulation of the cascade, so a reported
``optimality gap'' of greedy selection is always confounded with sampling
noise, and no exact ground truth has been available at any scale beyond
trivial instances.

\subsection*{Contributions}
\begin{enumerate}
\item \textbf{DEFER} (Section~\ref{sec:defer}), an adaptive wrapper around
  any one-shot IMIN planner. At each round it plans on the currently
  reachable region, protects the previous round's uncommitted plan against
  being replaced by a worse one, and commits a planned vertex only if a
  simple push-down inequality shows that waiting is not preferable. We
  prove: deferred commitment alone is free (Lemma~\ref{lem:defer-free});
  DEFER weakly dominates the one-shot planner it wraps
  (Theorem~\ref{thm:defer-dominates}); the push-down rule is the myopic
  optimum for the local commit-now/wait decision (Lemma~\ref{lem:pushdown});
  and the adaptivity gap of IMIN is unbounded, growing linearly in the
  maximum out-degree as the activation probability shrinks
  (Theorem~\ref{thm:adaptivity-gap}), with a matching, checkable
  construction.
\item \textbf{TWIG} (Section~\ref{sec:twig}), an exact dynamic program for
  the expected spread of a candidate blocking set on graphs of bounded
  treewidth, proved correct (Theorem~\ref{thm:twig-correct}) and then
  reformulated, via an exact linear-aggregation collapse
  (Theorem~\ref{thm:twig-collapse}), into a genuinely fixed-parameter
  tractable algorithm rather than merely an XP one
  (Propositions~\ref{prop:xp-space}--\ref{prop:fpt-space}). We further prove
  that the direction in which an elimination order is processed is not a
  tuning detail but a first-class algorithmic choice
  (Proposition~\ref{prop:direction}), with a worst-case family (the star
  graph) on which the wrong direction forces the frontier to scale with $n$
  even though the treewidth is $1$.
\item An exact, sampling-noise-free measurement of the optimality gap of
  greedy IMIN selection (Section~\ref{sec:twig-empirical}): on graphs small
  enough for TWIG to enumerate, greedy is exactly optimal in 56 of 60
  instances and departs from the true optimum by up to 84.6\% in the worst
  case found -- the first such measurement in the IMIN literature free of
  Monte-Carlo confounds.
\item A benchmark of DEFER against AdvancedGreedy, GreedyReplace, and
  simple heuristics on SNAP-scale contact and collaboration networks
  (Section~\ref{sec:experiments}), together with an auxiliary,
  apples-to-apples comparison on the companion source-detection problem
  used throughout our broader project.
\end{enumerate}

\section{Related Work}\label{sec:related}

\textbf{Influence minimization.} \citet{xie2023} introduce IMIN, prove
NP-hardness and inapproximability, and give AdvancedGreedy. \citet{xie2024}
extend this with GreedyReplace and Sandwich-style bounds and give the
network-reliability-flavored complexity analysis that our Theorem
\ref{thm:twig-correct} builds on for exact verification; their own
Exact-vs-GreedyReplace check (their Table 3) uses a borrowed exact method on
100-node subgraphs at up to four budget levels, considerably smaller and
less systematic than the fully proof-verified, 60-instance study we report
in Section~\ref{sec:twig-empirical}.

\textbf{Exact spread computation.} Computing expected spread exactly is a
network-reliability problem and is \#P-hard in general
\citep{valiant1979,provan1983}. \citet{nakamura2026} give a linear-time
exact algorithm for expected influence spread on bounded-\emph{pathwidth}
graphs; their result computes spread only and does not address blocking set
selection, greedy verification, or the adaptive setting, and pathwidth is a
strictly more restrictive parameter than the treewidth TWIG uses. Our
treewidth dynamic program follows the classical bounded-treewidth
programming paradigm \citep{bodlaender1988,cygan2015}.

\textbf{Adaptive submodular and sequential decision problems.} Adaptive
variants of influence \emph{maximization} and related coverage problems are
well studied; to our knowledge no prior work formalizes an adaptive variant
of IMIN or proves a lemma of the shape of Lemma~\ref{lem:defer-free} (that
deferred commitment of a fixed one-shot plan is loss-free) or an
unbounded-adaptivity-gap result of the shape of
Theorem~\ref{thm:adaptivity-gap} for this problem.

\section{Preliminaries}\label{sec:prelim}

\textbf{Independent cascade model.} $G=(V,E)$ is a directed graph,
$n=|V|$, $m=|E|$, with edge probabilities $p:E\to(0,1]$. A live-edge
realization $\omega\in\{0,1\}^E$ is drawn with
$\Pr[\omega]=\prod_{e}p_e^{\omega_e}(1-p_e)^{1-\omega_e}$. For a seed set
$S\subseteq V$ and blocked set $B\subseteq V\setminus S$, write
$G_B=G[V\setminus B]$ and let $\Rspread(\omega,B)\subseteq V\setminus B$ be
the set of vertices reachable from $S$ using only live edges of $G_B$. The
\emph{expected spread} is
\begin{equation}\label{eq:spread}
\FF(B)\;=\;\mathbb E_\omega\bigl[\,|\Rspread(\omega,B)\setminus S|\,\bigr].
\end{equation}

\textbf{IMIN (one-shot).} Given $G$, $p$, $S$, and budget $b$, choose
$B\subseteq V\setminus S$, $|B|\le b$, minimizing $\FF(B)$.

\textbf{Adaptive IMIN.} The cascade unfolds in rounds $t=0,1,2,\dots$. In
round $t$ the controller observes the newly activated set $N_t$
(full-adoption feedback: activations, not live edges, are observed), may
block any inactive, unblocked vertex, and the total number of vertices
blocked over all rounds is at most $b$. The objective is to minimize the
expected final spread. A one-shot solution is the special case that spends
the whole budget in round $0$.

\textbf{Treewidth machinery (used only in Section~\ref{sec:twig}).} Fix a
permutation $\sigma=(\sigma_1,\dots,\sigma_n)$ of $V$ and write
$\mathrm{pos}(v)$ for $v$'s index. For $v=\sigma_i$, let
$\mathrm{connect\_to}(v)=\{u\in N(v):\mathrm{pos}(u)<i\}$ and
$\mathrm{last}(v)=\max_{x\in N[v]}\mathrm{pos}(x)$. The \emph{frontier}
after step $i$ is $F_i=\{\sigma_j:j\le i,\ \mathrm{last}(\sigma_j)\ge i\}$,
and the \emph{width} of $\sigma$ is $w=\max_i|F_i|-1$, which upper-bounds
the treewidth of $G$ when $\sigma$ comes from a min-fill-in / min-degree
chordal completion \citep{bodlaender1988,cygan2015}.

\section{DEFER: Adaptive Influence Minimization by Deferred Frontier Blocking}\label{sec:defer}

\subsection{Algorithm}\label{sec:defer-algo}

DEFER wraps a one-shot planner $\mathcal P$ (any IMIN algorithm; by default
a dominator-tree greedy planner with isolation moves). At each round it
re-plans on the currently reachable region, keeps whichever of the new plan
and the previous round's uncommitted leftover looks better on freshly drawn
samples (\emph{plan protection}), and commits only the fraction of the plan
that is both exposed to the current frontier and passes a push-down test.

\begin{algorithm}[h]
\caption{DEFER}\label{alg:defer}
\begin{algorithmic}[1]
\Require graph $G$, probabilities $p$, seeds $S$, budget $b$, sample size
  $\theta$, one-shot planner $\mathcal P$
\State $A_0\gets S$, $B_0\gets\varnothing$, $F_0\gets S$, $L\gets\varnothing$ (uncommitted leftover)
\For{$t=0,1,2,\dots$ while the cascade is alive and $|B_t|<b$}
  \State $b_t\gets b-|B_t|$
  \State draw $\theta$ live-edge subgraphs reachable from $F_t$ avoiding $A_t\cup B_t$
  \State $P'\gets\mathcal P(\text{seeds}=F_t,\ \text{forbidden}=A_t\cup B_t,\ \text{budget}=b_t)$ on the samples
  \State $P_t\gets\arg\max\{\text{sample-average saving of } P', L\}$ \Comment{plan protection}
  \State $X_t\gets$ inactive, unblocked out-neighbors of $F_t$
  \For{$v\in P_t\cap X_t$}
    \State $q_v\gets 1-\prod_{u\in F_t}(1-p(u,v))$; $c_v\gets$ inactive unblocked out-degree of $v$
    \State $\lambda_t\gets$ smallest marginal saving among the picks of $P_t$ (0 if $P_t$ does not exhaust $b_t$)
    \If{$q_v\,(c_v+1/\lambda_t)\ge 1$} \Comment{push-down rule, Lemma~\ref{lem:pushdown}}
      \State block $v$ now
    \EndIf
  \EndFor
  \State $L\gets P_t\setminus(\text{blocked this round})$
  \State observe $N_{t+1}$; update $A_{t+1},F_{t+1}$
\EndFor
\end{algorithmic}
\end{algorithm}

The push-down rule compares spending one budget unit on an exposed node $v$
now against waiting one round and, if $v$ activates (probability $q_v$),
protecting its $c_v$ children instead while losing $v$ itself. When the
planner is deterministic given the samples and no new information arrives,
DEFER reduces to lazily committing the one-shot planner's own plan
(Lemma~\ref{lem:defer-free}); the two ablations used in
Section~\ref{sec:experiments} are COMMIT (same loop, but the whole plan is
committed every round) and the one-shot planners run without any wrapper.

\subsection{Theoretical guarantees}\label{sec:defer-theory}

Full proofs of every result in this section are given in
Online Supplement~\ref{ec:defer-proofs}.

\begin{lemma}[Deferral is free]\label{lem:defer-free}
Fix a live-edge realization, the seeds, and any set $P$ with $|P|\le b$. Let
$\pi_P$ be the policy that, in every round, blocks the not-yet-blocked
vertices of $P$ that are exposed to the current frontier. Then $\pi_P$
produces exactly the same final active set as blocking $P$ at time $0$, and
blocks at most as many vertices.
\end{lemma}

\begin{theorem}[DEFER weakly dominates its planner]\label{thm:defer-dominates}
Let $P_0$ be $\mathcal P$'s plan at round $0$. If, in every later round, the
planner returns a plan whose sample-average saving is at least that of the
leftover $P_{t-1}\setminus C_{t-1}$ -- guaranteed when the leftover is
offered to the planner as a candidate solution -- then the expected final
spread of DEFER is at most that of the one-shot solution $P_0$, up to the
sampling error of the $\theta$ samples.
\end{theorem}

\begin{lemma}[Push-down is the myopic optimum]\label{lem:pushdown}
Consider an exposed planned vertex $v$ whose children (inactive
out-neighbors) jointly protect the same set as $v$ minus $v$ itself, assume
the children are exposed only through $v$, and value a budget unit at
$\lambda$ vertices. Among the two policies ``block $v$ now'' and ``block
$v$'s children if and only if $v$ activates,'' the second has smaller
expected loss (vertices lost plus $\lambda$ times budget spent) if and only
if $q_v(c_v+1/\lambda)<1$.
\end{lemma}

\begin{theorem}[The adaptivity gap of IMIN is unbounded]\label{thm:adaptivity-gap}
For every $\Delta\ge 2$ and $p\in(0,1)$ there is an instance with maximum
out-degree $\Delta$ on which the best one-shot solution with budget $1$ has
expected spread at least $\bigl(1-(1-p)^{\Delta}\bigr)/p$ times that of DEFER
with budget $1$ (up to an additive constant); this ratio tends to $\Delta$ as
$p\to 0$.
\end{theorem}

\begin{theorem}[Complexity of DEFER]\label{thm:defer-complexity}
Let $n_t,m_t$ be the number of vertices and edges reachable from $F_t$ in
the union of the $\theta$ samples, $r$ the number of rounds, and $\alpha$
the inverse Ackermann function. With the dominator planner, round $t$ costs
$O(\theta\,m_t\,\alpha(m_t,n_t))$ for the dominator trees plus
$O(b_t\,\theta'\,m_t\,\alpha)$ for the greedy picks ($\theta'\le\theta$);
the total is $O\bigl(\sum_t(1+b_t)\,\theta\,m_t\,\alpha\bigr)$, which is
$O(r\,b\,\theta\,m\,\alpha(m,n))$ in the worst case.
\end{theorem}

An additional ablation, \texttt{isocut}, adds frontier-isolation moves to
the planner and admits a further worst-case separation
(Theorem~\ref{thm:isocut-separation}, Online Supplement) from
AdvancedGreedy, GreedyReplace, and the Sandwich lower bound, though it is
reported only as an ablation because it is rarely selected on the SNAP
graphs used in Section~\ref{sec:experiments}.

\begin{theorem}[Isolation-move separation]\label{thm:isocut-separation}
There is a family of instances on which AdvancedGreedy, GreedyReplace, and
the Sandwich lower-bound greedy all save $O(b)$ vertices with budget $b$,
while a planner with frontier-isolation moves saves $\Omega(N)$ for
arbitrary $N$.
\end{theorem}

\section{TWIG: Exact Treewidth-Bounded Validation}\label{sec:twig}

Every quantity reported by a one-shot IMIN heuristic -- including the plans
DEFER wraps -- is normally a Monte-Carlo estimate of Equation
\eqref{eq:spread}, because computing $\FF(B)$ exactly is a network-reliability
computation and is \#P-hard in general \citep{valiant1979,provan1983}. TWIG
computes $\FF(B)$ \emph{exactly} on graphs of bounded treewidth, giving a
ground-truth oracle that we use to certify how close greedy selection comes
to the true optimum, free of sampling noise.

\subsection{Algorithm}\label{sec:twig-algo}

\begin{definition}[Configuration]\label{def:config}
A configuration over a vertex set $F\subseteq V$ is a pair
$\kappa=(\Pi,\rho)$ where $\Pi$ is a set partition of $F$ with at most one
distinguished block $\mathrm{SEED}$, and $\rho$ maps each non-$\mathrm{SEED}$
block to a nonnegative integer (its \emph{pending count}).
\end{definition}

\texttt{ExactSpread} maintains, after each step $i$ of the schedule defined
in Section~\ref{sec:prelim}, a probability distribution over configurations
on the frontier $F_i$, plus a running scalar $\mathrm{acc}_i$, updated by
three operations: \emph{introduce} a new singleton block for each newly
processed vertex; \emph{process\_edge} branches on each edge's live/dead
outcome, merging blocks on the live branch and paying out to $\mathrm{acc}$
when a merge is into $\mathrm{SEED}$; and \emph{retire} pays out $\mathrm{acc}$
for a retiring $\mathrm{SEED}$-block vertex (unless it is itself a seed), or
defers a retiring ordinary vertex into its block's pending count, discarding
that count once the block empties of active members.

\begin{theorem}[Correctness of \texttt{ExactSpread}]\label{thm:twig-correct}
For every $B$, \texttt{ExactSpread}$(G,p,S)$.evaluate$(B)=\FF(B)$.
\end{theorem}

\subsection{Complexity: from XP to FPT in treewidth}\label{sec:twig-complexity}

\begin{proposition}[State space of \texttt{ExactSpread}]\label{prop:xp-space}
The number of distinct configurations at any step is at most
$B_{w+1}\cdot(n+1)^{w+1}$, where $B_{w+1}$ is the $(w+1)$-th Bell number.
\end{proposition}

This makes \texttt{ExactSpread} XP, not FPT, in treewidth: the degree of the
polynomial in $n$ grows with $w$.

\begin{theorem}[Correctness of the linear-aggregation collapse]\label{thm:twig-collapse}
Fix a partition shape $s$ and, at some step, let
$\mathcal K_s=\{(s,\rho):P_{(s,\rho)}>0\}$. Define the aggregates
$P_s=\sum_{(s,\rho)\in\mathcal K_s}P_{(s,\rho)}$ and
$W_b=\sum_{(s,\rho)\in\mathcal K_s}P_{(s,\rho)}\cdot\rho(b)$ for each
non-$\mathrm{SEED}$ block $b$ of $s$. Then $(P_s,\{W_b\})_s$ can be updated
at every step by the four rules \texttt{ExactSpreadFast} implements, and
this update yields the identical sequence of $\mathrm{acc}$ increments as
\texttt{ExactSpread}, for every $B$.
\end{theorem}

\begin{proposition}[State space of \texttt{ExactSpreadFast}]\label{prop:fpt-space}
The number of distinct states at any step is at most $B_{w+1}$, independent
of $n$, giving overall running time $O(n\cdot B_{w+1}\cdot\mathrm{poly}(w))$.
\end{proposition}

Proposition~\ref{prop:fpt-space} is a genuine change of complexity class
from Proposition~\ref{prop:xp-space}, not merely a constant-factor speed-up:
the $n^{w+1}$ factor is eliminated entirely.

\subsection{Processing direction is a first-class algorithmic choice}\label{sec:twig-direction}

\begin{proposition}[Star graphs separate forward and reversed elimination order]\label{prop:direction}
Let $T_k$ be the star $K_{1,k}$, $n=k+1$. Min-fill-in computes the
elimination order $\pi=(\ell_1,\dots,\ell_k,c)$, correctly certifying
treewidth $1$. Running the schedule of Section~\ref{sec:prelim} directly on
$\pi$ (the centre processed last) produces a frontier of size $k=n-1$
immediately before the centre's step. Running it on $\mathrm{reverse}(\pi)$
(the centre processed first) keeps the frontier at size at most $2$
throughout.
\end{proposition}

This is a worst-case separation on an exact, checkable family: a
bounded-treewidth graph (here, treewidth $1$, for any $n$) on which the
\emph{same} elimination order, run in the wrong direction, forces the true
state space of the dynamic program to scale with $n$ instead of remaining
bounded by a function of the width alone. In our implementation the
elimination order is therefore always reversed before scheduling.

\subsection{Empirical validation}\label{sec:twig-empirical}

Combining the collapse of Theorem~\ref{thm:twig-collapse} with the reversed
processing order of Proposition~\ref{prop:direction} makes exact evaluation
tractable on real, non-synthetic small networks:

\begin{table}[h]
\TABLE
{Exact evaluation feasibility on real small networks\label{tab:twig-feasibility}}
{\begin{tabular}{@{}lccc@{}}
\hline
\up Network & $n$ & width (min-fill-in) & \texttt{ExactSpreadFast.evaluate()} \\
\hline
\up Zachary karate club & 34 & 5 & 2.5 s \\
Iceland (sexual contact) & 75 & 4 & 51 s \\
\down Dolphin social network & 62 & $\sim$10--11 & infeasible (memory) \\
\hline
\end{tabular}}
{Before both fixes (Theorem~\ref{thm:twig-collapse} and reversed processing order), karate did not finish in 60 s and Iceland/Dolphin were not attempted.}
\end{table}

On the karate club network (one seed, random independent-cascade
probabilities in $[0.05,0.4]$), \texttt{greedy\_exact} -- greedy selection
using \texttt{ExactSpreadFast} itself as the marginal-gain oracle, with zero
sampling noise -- selects node 2 at budget 1 (expected spread falls from the
unblocked value to 7.98) and adds node 1 at budget 2 (expected spread
5.51). \citet{xie2023} prove $\FF(\cdot)$ is not supermodular in $B$ under
the independent-cascade model, so greedy carries no $(1-1/e)$-style
guarantee for IMIN, and \citet{xie2024} further show IMIN is APX-hard;
consistent with both results, comparing \texttt{greedy\_exact} against
\texttt{brute\_force\_optimal\_block} on 60 small random instances
($n=6$--$14$, budget 1--4) finds greedy exactly optimal in 56 of 60 and
strictly suboptimal in 4, by up to 84.6\% relative excess spread in the
worst case (mean excess 2.2\% across all 60). This is, to our knowledge, the
first exact (not Monte-Carlo-confounded) measurement of how often and how
badly greedy departs from optimal IMIN selection.

\section{Experiments}\label{sec:experiments}

\subsection{DEFER benchmark}\label{sec:defer-experiments}

We benchmark DEFER, wrapping the dominator-tree planner (default) and
GreedyReplace (alternate planner), against the one-shot planners
themselves, the COMMIT ablation, and simple heuristics (proximity, degree),
on ca-GrQc, ca-HepTh, a power-grid network, a Barab\'asi--Albert graph, and
soc-Epinions1 where feasible, under constant activation probabilities
$p\in\{0.05,0.1\}$ and the weighted-cascade model, seed sets of size 20 or
50 (random or highest-degree), and budgets 20, 50, 100. \emph{[Full results
table pending completion of the benchmark run; see
\texttt{docs/RESULTS\_NOTES.md} in the accompanying code release for the
current partial results and Section~\ref{ec:reproducibility} of the Online
Supplement for the exact protocol.]}

\subsection{Auxiliary benchmark: source detection}\label{sec:source-experiments}

As a secondary, apples-to-apples benchmark, we also evaluate classical and
learned source-detection methods (rumor centrality, NetSleuth, dynamic
message passing, and the graph neural network architectures of
\citealp{sterchi2026} and the works it reproduces) on identical,
seed-fixed test instances across several real contact and collaboration
networks, calibrated per network to the target infected fraction of
\citet{sterchi2026}. \emph{[Full six-network table pending completion of the
GNN training matrix and the reduced-sample DMP run on the largest network;
partial results are in \texttt{results/gnn\_benchmark\_sterchi/} in the code
release.]}

\section{Conclusion}\label{sec:conclusion}

We gave the first formal treatment of adaptivity for influence
minimization: a wrapper, DEFER, that is provably never worse than the
one-shot planner it wraps, a myopic-optimal commitment rule, and a proof
that the benefit of adapting to the unfolding cascade is unbounded rather
than a modest constant-factor improvement. We complemented this with TWIG,
an exact, fixed-parameter-tractable dynamic program for verifying any IMIN
heuristic's optimality gap without the sampling confounds that have affected
every prior empirical evaluation in this literature, together with a
provable separation showing that even the direction of processing a fixed
tree decomposition is algorithmically consequential. Both contributions are
implemented, tested against brute force, and released with the paper.

% Endnotes (none currently used)
%\begingroup \parindent 0pt \parskip 0.0ex \def\enotesize{\normalsize} \theendnotes \endgroup

\ACKNOWLEDGMENT{The author thanks no one in particular for this draft; acknowledgments will be added prior to submission.}

\ECSwitch
\ECDisclaimer
\ECHead{Online Supplement}

\section{Proofs for Section~\ref{sec:defer} (DEFER)}\label{ec:defer-proofs}

\begin{repeatlemma}[Lemma~\ref{lem:defer-free} (Deferral is free)]
Fix a live-edge realization, the seeds, and any set $P$ with $|P|\le b$. Let
$\pi_P$ be the policy that, in every round, blocks the not-yet-blocked
vertices of $P$ that are exposed to the current frontier. Then $\pi_P$
produces exactly the same final active set as blocking $P$ at time $0$, and
blocks at most as many vertices.
\end{repeatlemma}
\begin{proof}{Proof}
On a fixed realization a vertex $v$ becomes active in the first round in
which a frontier vertex has a live edge into it, i.e.\ the round after it is
first exposed. $\pi_P$ blocks $v\in P$ in that round, before its trial, so
$v$ never activates; every vertex outside $P$ faces the same set of blocked
vertices at each of its trials under both policies, because a vertex of $P$
that is never exposed never takes part in any trial. Hence the activation
sequences coincide. Vertices of $P$ that are never exposed are never
blocked.\Halmos
\end{proof}

\begin{repeattheorem}[Theorem~\ref{thm:defer-dominates} (DEFER weakly dominates its planner)]
Let $P_0$ be the plan of $\mathcal P$ at round $0$. If, in every later
round, the planner returns a plan whose sample-average saving is at least
that of the leftover $P_{t-1}\setminus C_{t-1}$ -- guaranteed when the
leftover is offered to the planner as a candidate solution -- then the
expected final spread of DEFER is at most that of the one-shot solution
$P_0$, up to the sampling error of the $\theta$ samples.
\end{repeattheorem}
\begin{proof}{Proof}
By Lemma~\ref{lem:defer-free} the policy that keeps $P_0$ and commits
lazily has the same spread as blocking $P_0$ at time $0$. DEFER differs
from it only by replacing, at each round, the uncommitted leftover by a
plan that is no worse on the samples for the remaining budget; the
committed prefix is identical. Induction over rounds gives the claim on the
samples; concentration of the sample average (\citealp{xie2024}, applied
per candidate plan) gives the expectation statement.\Halmos
\end{proof}

\begin{repeatlemma}[Lemma~\ref{lem:pushdown} (Push-down is the myopic optimum)]
Consider an exposed planned vertex $v$ whose children (its inactive
out-neighbors) jointly protect the same set as $v$ minus $v$ itself, assume
the children are exposed only through $v$, and value a budget unit at
$\lambda$ vertices. Among the two policies ``block $v$ now'' and ``block
$v$'s children if and only if $v$ activates,'' the second has smaller
expected loss (vertices lost plus $\lambda$ times budget spent) if and only
if $q_v(c_v+1/\lambda)<1$.
\end{repeatlemma}
\begin{proof}{Proof}
Blocking now costs $\lambda$ and loses nothing of $v$'s set. Waiting loses
$v$ (one vertex) and spends $c_v$ units with probability $q_v$, and nothing
otherwise, while the protected set is the same in both branches by
assumption: expected loss $q_v+\lambda q_v c_v$. Comparing $\lambda$ against
$q_v+\lambda q_v c_v$ gives the stated condition.\Halmos
\end{proof}

\begin{repeattheorem}[Theorem~\ref{thm:adaptivity-gap} (The adaptivity gap of IMIN is unbounded)]
For every $\Delta\ge 2$ and $p\in(0,1)$ there is an instance with maximum
out-degree $\Delta$ on which the best one-shot solution with budget $1$ has
expected spread at least $\bigl(1-(1-p)^{\Delta}\bigr)/p$ times that of
DEFER with budget $1$ (up to an additive constant); this ratio tends to
$\Delta$ as $p\to 0$.
\end{repeattheorem}
\begin{proof}{Proof}
Seed $s$ has out-neighbors $a_1,\dots,a_\Delta$ with $p(s,a_i)=p$; each
$a_i$ starts a private directed path of length $L$ with edge probability
$1$. Any one-shot solution with budget $1$ blocks one vertex; blocking
$a_i$ (or a vertex on its path) saves at most $L+1$ vertices and only if
$a_i$ is activated, which happens with probability $p$: expected spread
$\ge(\Delta p-p)L$. DEFER observes $N_1=\{a_i:(s,a_i)\text{ live}\}$ and
blocks the first path vertex of one activated branch; if $k$ branches are
active the spread is $(k-1)L+k$. Its expected spread is
$\mathbb E[(k-1)^+]L+\Delta p\le(\Delta p-1+(1-p)^\Delta)L+\Delta p$ (DEFER's
plan at round $0$ is a branch head $a_i$ with $q=p$, $c=1$, and
$\lambda\approx pL$, so the push-down rule $p(1+1/(pL))<1$ defers it). The
ratio of one-shot to DEFER is at least
$\frac{(\Delta-1)p}{\Delta p-1+(1-p)^\Delta}$, which for $p\to0$ with
$\Delta p=c$ fixed tends to $c/(c-1+e^{-c})$ and, for fixed $\Delta$ as
$p\to0$, behaves like $\frac{(\Delta-1)p}{\binom{\Delta}{2}p^2}=
\frac{2}{\Delta p}\to\infty$; the stated bound
$(1-(1-p)^\Delta)/p$ follows from comparing the saving $Lp$ of the best
one-shot vertex with the saving $L\Pr[k\ge1]=L(1-(1-p)^\Delta)$ of
DEFER.\Halmos
\end{proof}

\begin{repeattheorem}[Theorem~\ref{thm:defer-complexity} (Complexity of DEFER)]
Let $n_t,m_t$ be the number of vertices and edges reachable from $F_t$ in
the union of the $\theta$ samples, $r$ the number of rounds, and $\alpha$
the inverse Ackermann function. With the dominator planner, round $t$
costs $O(\theta\,m_t\,\alpha(m_t,n_t))$ for the dominator trees plus
$O(b_t\,\theta'\,m_t\,\alpha)$ for the greedy picks; the total is
$O\bigl(\sum_t(1+b_t)\,\theta\,m_t\,\alpha\bigr)$, which is
$O(r\,b\,\theta\,m\,\alpha(m,n))$ in the worst case.
\end{repeattheorem}
\begin{proof}{Proof}
Each dominator tree is computed by the Cooper--Harvey--Kennedy iteration
on the reachable subgraph (near-linear; the Lengauer--Tarjan analysis gives
the $\alpha$ bound). Step 2 of Algorithm~\ref{alg:defer} is a breadth-first
search per sample. Recomputation after a pick is restricted to samples
whose cached reach set contains the picked vertex, which is exactly the
set of samples whose dominator tree can change. Summing over picks and
rounds gives the bound.\Halmos
\end{proof}

\begin{repeattheorem}[Theorem~\ref{thm:isocut-separation} (Isolation-move separation)]
There is a family of instances on which AdvancedGreedy, GreedyReplace, and
the Sandwich lower-bound greedy all save $O(b)$ vertices with budget $b$,
while a planner with frontier-isolation moves saves $\Omega(N)$ for
arbitrary $N$.
\end{repeattheorem}
\begin{proof}{Proof}
Construction: a decoy seed with $b+1$ out-neighbors, each with one private
leaf (single-vertex saving 2); a real seed with $c\le b$ out-neighbors, all
of which point to every vertex of a region $R$ of $N$ vertices whose
vertices have no private successors. Every single vertex in $R$ or among
the real seed's out-neighbors saves at most 1, so all three published
rules spend the budget on decoy neighbors (saving $2b$; GreedyReplace's
replacement pass replaces a decoy blocker by itself and stops), whereas
isolating the real seed saves $N+c$ at cost $c$.\Halmos
\end{proof}

\section{Proofs for Section~\ref{sec:twig} (TWIG)}\label{ec:twig-proofs}

\begin{repeatlemma}[Invariant maintained by \texttt{ExactSpread}]\label{ec:lem:invariant}
Fix $B$ and a step $i$. For every configuration $\kappa=(\Pi,\rho)$ with
$P_\kappa>0$ after step $i$,
\begin{align*}
P_\kappa=\Pr\nolimits_\omega\Bigl[\;&\Pi\text{ is the partition of }F_i\text{
by connectivity in }\bigl(F_i^{\ge},\,\omega|_{E_i}\bigr)\text{, seed-
containing part collapsed to SEED};\\
&\text{and }\forall\text{ non-SEED block }b,\ \rho(b)=\bigl|
\{u:\mathrm{last}(u)\le i,\ u\text{ retired non-seed},\ u\sim_\omega b\}
\bigr|\;\Bigr],
\end{align*}
where $E_i$ is the set of edges processed by step $i$ (both endpoints
introduced, neither endpoint in $B$), $F_i^{\ge}=\{x\notin B:\mathrm{pos}(x)
\le i\}$ is every vertex introduced so far (active or already retired), and
$u\sim_\omega b$ means $u$ is connected, via live edges of $E_i$ within
$F_i^{\ge}$, to the current representative of block $b$. Moreover
$$\mathrm{acc}_i=\sum_{v\in V\setminus S\setminus B\,:\,\mathrm{last}(v)\le
i}\Pr\nolimits_\omega\bigl[v\sim_\omega\text{ some seed, using only edges of
}E_i\bigr].$$
\end{repeatlemma}
\begin{proof}{Proof}
By induction on $i$. \emph{Base case} $i=-1$ (before any step): $F_{-1}
=\varnothing$, the unique configuration is the empty partition with $P=1$,
$\mathrm{acc}=0$, and both sides of the claimed identities are vacuously
true ($E_{-1}=\varnothing$, no vertex has retired).

\emph{Introduce.} Adding $v$ as a new singleton block changes nothing about
already-active vertices' connectivity (no new edge is realized), and
correctly starts $v$'s own block (SEED if $v\in S$, pending $0$ otherwise,
since $v$ is trivially connected to itself and, being newly introduced, to
nothing else yet). The invariant is preserved for $\Pi$ and, since no new
edges enter $E_i$, for $\mathrm{acc}$.

\emph{Process\_edge$(u,v)$.} This is the only place a new edge $(u,v)$
enters $E_i$. Conditioning on its two outcomes (independent of every other
edge's realization, by the independent-cascade model's edge independence)
exactly reproduces the two branches: tails leaves $\Pi$, $\rho$,
$\mathrm{acc}$ untouched, matching that $(u,v)\notin\omega$ leaves
connectivity and every retired vertex's membership unchanged. Heads merges
$u$'s and $v$'s blocks, which is exactly the effect of adding a live edge to
a union-find-style connectivity structure. If the merge is into SEED: every
already-retired vertex $u'$ counted in the merging block's $\rho(b)$ is, by
the outer induction hypothesis, connected (via $E_i\setminus\{(u,v)\}$) to
$b$'s current representative, hence now (with $(u,v)$ live) transitively
connected to a seed for the first time in this history -- exactly one unit
each is owed to $\mathrm{acc}$, contributing $p(u,v)\cdot\rho(b)$ in
expectation over this branch, matching the update; $\rho(b)$ is then
correctly discharged (those vertices will never need to be paid for again,
since a vertex is reached at most once). If the merge is between two
ordinary blocks, no vertex's fate is decided yet (neither is connected to a
seed), so $\mathrm{acc}$ is untouched and the two blocks' pending vertex
sets simply become one, i.e.\ $\rho(b_1)+\rho(b_2)$, matching the update.

\emph{Retire$(u)$.} $u$ retiring means $\mathrm{last}(u)=i$: every edge
incident to $u$ is already in $E_i$, so $u$'s connectivity to everything
else active or retired is already fully determined and will never change.
If $u$'s block is already SEED: $u$ is reached with probability $P_\kappa$
in this history, contributed to $\mathrm{acc}$ unless $u$ itself is a seed
(seeds are excluded from ``reached,'' matching the definition of
$\FF(B)$). Otherwise $u$'s fate is not yet decided; deferring it into
$\rho(b)$ is exactly the claimed count (one more retired vertex now
connected to $b$), \emph{unless} $b$ has no other active member, in which
case $b$ can never merge with anything else in the future (its only
remaining edges were already all processed, by definition of retirement
having emptied it of active representatives) -- correctly discharging
vertices that are provably never reached. Removing $u$ from $\Pi$ matches
$F_i\setminus\{u\}$.\Halmos
\end{proof}

\begin{repeattheorem}[Theorem~\ref{thm:twig-correct} (Correctness of \texttt{ExactSpread})]
For every $B$, \texttt{ExactSpread}$(G,p,S)$.evaluate$(B)=\FF(B)$.
\end{repeattheorem}
\begin{proof}{Proof}
Apply Lemma~\ref{ec:lem:invariant} at $i=n-1$ (the last step): $F_{n-1}
=\varnothing$ (every vertex has retired by the time all its neighbors,
which include itself, have been introduced), $E_{n-1}=E\setminus(B$-incident
edges$)$, i.e.\ all of $G_B$'s edges, and every non-seed, non-blocked
vertex has $\mathrm{last}(v)\le n-1$. The second identity of
Lemma~\ref{ec:lem:invariant} then reads $\mathrm{acc}_{n-1}=
\sum_{v\in V\setminus S\setminus B}\Pr_\omega[v\sim_\omega\text{ some seed
in }G_B]=\FF(B)$, which is exactly what \texttt{evaluate}$(B)$
returns.\Halmos
\end{proof}

\begin{repeatproposition}[Proposition~\ref{prop:xp-space} (State space of \texttt{ExactSpread})]
The number of partition shapes of a frontier of size at most $w+1$ is
$B_{w+1}$, the $(w+1)$-th Bell number, independent of $n$. But each shape's
pending vector has up to $w+1$ components, each ranging over
$\{0,\dots,n\}$, so the number of distinct configurations at any step is at
most $B_{w+1}\cdot(n+1)^{w+1}$.
\end{repeatproposition}

\begin{repeattheorem}[Theorem~\ref{thm:twig-collapse} (Correctness of the linear-aggregation collapse)]
Fix a partition shape $s$ and, at some step, let
$\mathcal K_s=\{(s,\rho):P_{(s,\rho)}>0\}$. Define the aggregates
$P_s=\sum_{(s,\rho)\in\mathcal K_s}P_{(s,\rho)}$ and
$W_b=\sum_{(s,\rho)\in\mathcal K_s}P_{(s,\rho)}\cdot\rho(b)$ for each
non-SEED block $b$ of $s$. Then $(P_s,\{W_b\})_s$ can be updated, at every
step, by the same four rules \texttt{ExactSpreadFast} implements, and this
update yields the identical sequence of \texttt{acc} increments as
\texttt{ExactSpread}, for every $B$.
\end{repeattheorem}
\begin{proof}{Proof}
We check each of the (at most) four ways a step transforms a configuration
$(s,\rho)\mapsto(s',\rho')$ within a fixed edge-outcome branch (shape
transitions never depend on $\rho$, only on $s$ and, for
\texttt{process\_edge}, the branch); linearity of each in $\rho$ is then
all that is needed.

\emph{Introduce a new non-seed block $b_0$.} $\rho'(b_0)=0$,
$\rho'(b)=\rho(b)$ otherwise, $s'=s\cup\{b_0\}$ is independent of $\rho$.
Then $W'_{b_0}=\sum_\rho P_\rho\cdot 0=0$ and $W'_b=W_b$ for $b\ne b_0$: no
update needed for existing entries, new entry is (implicitly) $0$.

\emph{Process\_edge, tails.} $s'=s$, $\rho'=\rho$, but the whole branch
carries probability weight $(1-p_e)$: $P'_s=(1-p_e)P_s$ and
$W'_b=(1-p_e)W_b$ for every entry.

\emph{Process\_edge, heads, merge into SEED.} The payout added to
$\mathrm{acc}$ in this branch is
$p_e\sum_{(s,\rho)}P_\rho\cdot\rho(b)=p_e\cdot W_b$ exactly (a direct
algebraic identity, not a mean-field substitution). Block $b$ is then
dropped ($\rho'$ has no entry for it).

\emph{Process\_edge, heads, merge two ordinary blocks $b_1,b_2\to b$.}
$\rho'(b)=\rho(b_1)+\rho(b_2)$, so
$W'_b=\sum_\rho P_\rho(\rho(b_1)+\rho(b_2))=W_{b_1}+W_{b_2}$.

\emph{Retire, block still active.} $\rho'(b)=\rho(b)+1$ for the retiring
block, others unchanged: $W'_b=\sum_\rho P_\rho(\rho(b)+1)=W_b+P_s$ (using
the aggregate $P_s$, not a per-history increment).

\emph{Retire, block emptied.} $b$'s entry is dropped for every history in
$\mathcal K_s$ simultaneously (the condition ``no other active member''
depends only on $s$, not $\rho$), so $W'_b$ is simply omitted.

\emph{Retire, block is SEED, $u\notin S$.} $\mathrm{acc}$ gains
$\sum_\rho P_\rho=P_s$ in this branch.

Each case is an exact algebraic identity between the aggregate before and
the aggregate after, with no term dropped or approximated; by induction
over the sequence of steps, $(P_s,\{W_b\})$ computed by
\texttt{ExactSpreadFast} therefore equals the shape-lumped aggregate of
\texttt{ExactSpread}'s full distribution at every step, and in particular
produces the identical sequence of contributions to $\mathrm{acc}$.\Halmos
\end{proof}

\begin{repeatproposition}[Proposition~\ref{prop:fpt-space} (State space of \texttt{ExactSpreadFast})]
The number of distinct states at any step is at most $B_{w+1}$ (one
$(P_s,\{W_b\})$ tuple per shape), independent of $n$.
\end{repeatproposition}

\begin{repeatproposition}[Proposition~\ref{prop:direction} (Star graphs separate forward and reversed elimination order)]
Let $T_k$ be the star $K_{1,k}$, $n=k+1$. Min-fill-in computes the
elimination order $\pi=(\ell_1,\dots,\ell_k,c)$, correctly certifying
treewidth $1$. Running the schedule of Section~\ref{sec:prelim} directly on
$\pi$ (the centre processed last) produces a frontier of size $k=n-1$
immediately before the centre's step. Running it on $\mathrm{reverse}(\pi)$
(the centre processed first) keeps the frontier at size at most $2$
throughout.
\end{repeatproposition}
\begin{proof}{Proof}
Immediate from the definitions of Section~\ref{sec:prelim}: under $\pi$,
$\mathrm{last}(\ell_i)=\mathrm{pos}(c)=k$ for every leaf (its only neighbor
$c$ is introduced last), so no leaf retires before step $k$, giving
$F_{k-1}=\{\ell_1,\dots,\ell_k\}$; under $\mathrm{reverse}(\pi)$,
$\mathrm{last}(\ell_i)=\mathrm{pos}(\ell_i)$ (once $c$ is already active,
$\ell_i$'s only edge is processed the instant $\ell_i$ is introduced, and
$c$ is not $\ell_i$'s own last neighbor to be introduced -- it already was),
so every leaf retires the step after it is introduced, and $c$ itself
retires only at the final step (its own $\mathrm{last}$ is the position of
the last-introduced leaf), giving $|F_i|\le 2$ for all $i$.\Halmos
\end{proof}

\section{Experimental Protocol}\label{ec:reproducibility}

\subsection{DEFER benchmark}
Graphs: ca-GrQc, ca-HepTh, a power-grid network, a Barab\'asi--Albert
random graph with 2000 vertices, and soc-Epinions1 where feasible.
Probabilities: constant $p\in\{0.05,0.1\}$ and the weighted-cascade model.
Seeds: 20 or 50, drawn uniformly at random or as the highest-degree
vertices. Budgets: 20, 50, 100. For each instance and each of 5 independent
live-edge realizations, every policy runs on the same realization with the
same total budget (common random numbers), so differences are attributable
to policy alone. Policies compared: none, AdvancedGreedy, GreedyReplace,
the Sandwich lower bound, proximity-to-seeds, degree (all one-shot); DEFER
with the dominator planner, DEFER with GreedyReplace as planner, COMMIT,
and DEFER restricted to a 4-round planning horizon. Reported: final spread
beyond the seeds (mean and standard error), fraction of spread saved,
budget used, and wall-clock time per episode including planning.

\subsection{TWIG optimality-gap study}
Random small graphs, $n=6$--$14$, edge density such that treewidth stays at
most 5--6, budget 1--4, one to four seeds. For each instance,
\texttt{greedy\_exact} and \texttt{brute\_force\_optimal\_block} are both
run with the identical exact oracle \texttt{ExactSpreadFast.evaluate}, so
any discrepancy between the two is a true optimality gap rather than
sampling noise. 60 instances were drawn; correctness of the oracle itself
was independently established by exhaustive summation over all $2^{|E|}$
live-edge subsets on graphs up to $n=9$ (2000+ configurations for
\texttt{ExactSpread}, 1000+ for \texttt{ExactSpreadFast}), matching
Theorems~\ref{thm:twig-correct} and \ref{thm:twig-collapse}.

\begin{thebibliography}{99}
\providecommand{\natexlab}[1]{#1}
\providecommand{\url}[1]{\texttt{#1}}
\providecommand{\urlprefix}{URL }

\bibitem[{Bodlaender(1988)}]{bodlaender1988}
Bodlaender HL (1988) Dynamic programming on graphs with bounded treewidth.
  \emph{Proc.\ 15th Internat.\ Colloquium on Automata, Languages and
  Programming (ICALP)}, 105--118.

\bibitem[{Budak et~al.(2011)Budak, Agrawal, and El~Abbadi}]{budak2011}
Budak C, Agrawal D, El~Abbadi A (2011) Limiting the spread of misinformation
  in social networks. \emph{Proc.\ 20th Internat.\ Conf.\ on World Wide Web
  (WWW)}, 665--674.

\bibitem[{Cygan et~al.(2015)Cygan, Fomin, Kowalik, Lokshtanov, Marx,
  Pilipczuk, Pilipczuk, and Saurabh}]{cygan2015}
Cygan M, Fomin FV, Kowalik {\L}, Lokshtanov D, Marx D, Pilipczuk M,
  Pilipczuk M, Saurabh S (2015) \emph{Parameterized Algorithms} (Springer,
  Cham).

\bibitem[{Kempe et~al.(2003)Kempe, Kleinberg, and Tardos}]{kempe2003}
Kempe D, Kleinberg J, Tardos {\'E} (2003) Maximizing the spread of influence
  through a social network. \emph{Proc.\ 9th ACM SIGKDD Internat.\ Conf.\ on
  Knowledge Discovery and Data Mining (KDD)}, 137--146.

\bibitem[{Nakamura and Nishino(2026)}]{nakamura2026}
Nakamura K, Nishino M (2026) Linear-time exact computation of influence
  spread on bounded-pathwidth graphs. \emph{39th Internat.\ Symp.\ on
  Algorithms and Computation / Scandinavian Symp.\ and Workshops on
  Algorithm Theory (SWAT)}, Dagstuhl LIPIcs. [Verify final proceedings
  citation prior to submission.]

\bibitem[{Provan and Ball(1983)}]{provan1983}
Provan JS, Ball MO (1983) The complexity of counting cuts and of computing
  the probability that a graph is connected. \emph{SIAM J.\ Comput.}
  12(4):777--788.

\bibitem[{Sterchi et~al.(2026)Sterchi, Brack, and Hilfiker}]{sterchi2026}
Sterchi M, Brack F, Hilfiker R (2026) A systematic benchmark of graph
  neural networks for epidemic source detection. \emph{arXiv:2512.20657}.

\bibitem[{Valiant(1979)}]{valiant1979}
Valiant LG (1979) The complexity of enumeration and reliability problems.
  \emph{SIAM J.\ Comput.} 8(3):410--421.

\bibitem[{Wang et~al.(2024)}]{wang2024sandwich}
Wang et~al. (2024) Sandwich approximation algorithms for influence
  minimization. [Placeholder: confirm full author list and venue prior to
  submission; referenced in this draft only for the Sandwich lower-bound
  baseline LSBM.]

\bibitem[{Xie et~al.(2023)Xie, Zhang, Wang, Lin, and Zhang}]{xie2023}
Xie J, Zhang Y, Wang Z, Lin X, Zhang W (2023) Minimizing the influence
  spread over a social network via node blocking. \emph{Proc.\ 39th IEEE
  Internat.\ Conf.\ on Data Engineering (ICDE)}, 3096--3108.

\bibitem[{Xie et~al.(2024)Xie, Zhang, Wang, Liu, Lin, and Zhang}]{xie2024}
Xie J, Zhang Y, Wang Z, Liu H, Lin X, Zhang W (2024) Influence minimization
  via blocking strategies. \emph{INFORMS J.\ Comput.} 37(6). Code:
  \texttt{github.com/INFORMSJoC/2024.0591}, arXiv:2312.17488.

\end{thebibliography}

\end{document}
